\documentclass[12pt]{article}

\usepackage{comment}
\usepackage[english]{babel}

\usepackage{amsmath}
\usepackage{amsfonts}
\usepackage{graphicx}
\usepackage[colorlinks=true, allcolors=blue]{hyperref}
\usepackage[autostyle, english = american]{csquotes}
\MakeOuterQuote{"}
\usepackage{caption}
\captionsetup{font={footnotesize, sf}, labelfont=bf}

\usepackage{multirow}
\usepackage[table]{xcolor}
\usepackage{booktabs}

\definecolor{delayA}{HTML}{b75347} % Lognormal
\definecolor{delayB}{HTML}{edc775} % Generalized-Gamma
\definecolor{delayC}{HTML}{94b594} % Dirichlet
\definecolor{epincast}{HTML}{94b594} % Dirichlet
\definecolor{nobbs}{HTML}{94b594} % Dirichlet

\title{Nowcasting from an epidemic's reporting-delay. A censored regression approach.}
\author{Rodrigo Zepeda-Tello, Jeffrey Shaman}
\date{\today{}}
\begin{document}
\maketitle

\section{Introduction}

%P1: Surveillance is cool
Effective epidemic surveillance is crucial for public health decision-making \cite{world2023future, shen2025progress}. Central to this effort is the continuous monitoring of incident cases across complementary systems such as laboratory-based diagnostic testing, clinician-reported case notifications, and vital statistics mortality reporting \cite{ibrahim2020epidemiologic}. Timely and reliable estimates of these indicators are essential for characterizing the trajectory of an outbreak and anticipating healthcare needs. Without accurate and up-to-date surveillance inputs, policy responses risk being delayed or misallocated \cite{shen2025progress, villela2022importance}.

%P2: Sometimes we don't see all the cases and people do weird things
Clinical and healthcare-system reporting constraints introduce delays between the time an infection occurs, the time it is detected, and the time it is formally recorded in surveillance systems. As a result, the number of cases reported on a given day rarely represents the true number of infections that occurred \cite{jajosky2004evaluation, swaan2018timeliness, marinovic2015quantifying}. Additional cases associated with earlier dates continue to appear as delayed reports accumulate, and the most recent incidence estimates typically require a stabilization period before they can be considered reliable \cite{center2004morbidity}. In practice, this often leads to analysts ignoring the most recent data until sufficient time has passed for backlogged reports to surface thus wasting valuable real-time information.

Nowcasting provides a solution to this challenge. Often described as "predicting the present" \cite{scott2014predicting}, nowcasting refers to the inference of the current state of affairs using historical and partially observed contemporary data \cite{wu2021nowcasting}. In the context of epidemiological time series, this usually corresponds to producing an estimate of current incidence (or a function of it such as the effective reproductive number, $R_t$). Originating in the meteorological literature, nowcasting methods have been applied across disciplines including earthquake prediction \cite{rundle2016nowcasting}, short-term weather assessment \cite{camporeale2019challenge}, and macroeconomic forecasting \cite{bok2018macroeconomic}. In epidemiology, nowcasting has been used extensively to infer real-time incidence during outbreaks of influenza \cite{he2022nowcasting, preis2014adaptive, spreco2018evaluation}, ebola \cite{tariq2025nowcast}, COVID-19 \cite{greene2021nowcasting, birrell2021real, albani2022nowcasting}, dengue \cite{codeco2018infodengue, mcgough2020nowcasting}, mpox \cite{charniga2024nowcasting, rohrer2025nowcasting}, among others \cite{mellor2025application, menkir2021nowcasting}.

Most epidemiological nowcasting approaches follow a common two-step structure: (1) modeling the reporting-delay distribution, and (2) modeling the underlying epidemic process. A range of methods has been proposed for these tasks, including random forests \cite{sahai2022machine}, regression-based approaches \cite{donker2011nowcasting, mcgough2020nowcasting, bastos2019modelling}, and Bayesian hierarchical models \cite{johnson2025baseline,gunther2021nowcasting, hohle2014bayesian}; the latter offering particular advantages early in an outbreak by incorporating prior knowledge. 

%{\color{red} The following might be better in Methods}

Many of these methods typically treat reporting-delays as a nuisance mechanism required to reconstruct latent incidence through a "reporting triangle" representation\cite{verrall1989state, mcgough2020nowcasting}. In contrast, we formulate reporting-delays as realizations of a stochastic process that can be modeled directly and jointly with the epidemic dynamics. This perspective yields two advantages. First, it leads to a likelihood that depends only on observed delays, avoiding the need to reconstruct the full event-report matrix which is normally used even when zero cases have been observed. Second, the inference on the reporting process itself, allows for the detection of anomalies and treatment of some missing values. As a result, delays become an additional source of information about the surveillance system.

%Many of these methods typically rely on an actuarial construction referred to as the "reporting triangle" , which requires estimating the number of cases that will eventually be reported for each potential report date. This quantities, however, are not strictly necessary for most nowcasting applications, in which solely the final number of cases by event date is of interest. Nonetheless, existing formulations typically incorporate the full reporting triangle into the likelihood, thereby introducing nuisance structure that does not contribute directly to the inferential target.

{\color{red} \textbf{[Stand-in paragraph]}:

Here, we introduce this alternative framework. Next, we show with real-life examples how reformulation simplifies the inferential task. Finally, we compare our estimation against current methods to show...
}

\section{Methods}


The main idea of the methodology is to use a censored data approach to infer the latent epidemic process through the delay observations. We assume that the unobserved epidemic process is described by $\{N_t\}_t$. At each moment in time, $t$, there is a reporting-delay process $\{D_t\}_t$ that governs the distribution of the times when each case will be observed. While the epidemic process governs \textit{how many} cases there will be, the reporting-delay process governs \textit{when} they will be reported.  As an example, consider that, overall, 12 cases happened on day 5 of the epidemic. From those 12 cases, we could have recorded $3$ on day $5$ (delay = 0), $7$ on day $6$ (delay = 1) and $2$ on day $9$ (delay = 4). This idea of the two processes is sketched in Figure \ref{fig:nowcastingidea}. 

% Subfigures 
\begin{comment}
\begin{figure}[!htb]%
    \centering
    \subfloat[]{{\includegraphics[width=0.5\textwidth]{images/nowcasting_idea.pdf} }}%
    \subfloat[]{{\includegraphics[width=0.5\textwidth]{images/delay_epidemic_process.pdf} }}%
    \caption{\footnotesize (a) reporting-delay-epidemic processes for an infectious disease. At time $t$ the (unobserved) epidemic process governs how many cases there will be while the reporting-delay process governs when they will be registered. (b) Classical reporting triangle representation of the reporting-delay and the epidemic processes. Notation $m_{t}^{d}$ represents the total number of cases of time $t$ reported with delay $d$.}%
    \label{fig:nowcastingidea}%
\end{figure}
\end{comment}

\begin{figure}[!htb]%
    \centering    \includegraphics[width=\linewidth]{images/diagram.pdf} 
    \caption{Reporting-delay-epidemic processes for an infectious disease. At time $t$ the (unobserved) epidemic process governs how many cases there will be while the reporting-delay process governs when they will be registered.}%
    \label{fig:nowcastingidea}%
\end{figure}

In the following sections, we provide  details regarding the proposed nowcasting methodology. First we show a simplified version modeling the reporting-delay process. Second, we add an additional type of censoring for when the reporting-delay is missing. Finally, we extend the model to allow for different examples of epidemic models. 


%THE IDEA IS SIMPLE: WE FOCUS ON A STOCHASTIC PROCESS ON THE DELAY (RELATE TO HOW POISSON AND EXPONENTIAL DELAYS)


\subsection{Reporting-delay process}

We consider for time $t \geq 0$ a collection of integer delays $d_1^t, d_2^t, \dots, d_{k_t}^t$ bounded above by $d_*^{t}$. At time $t$, each $d_{\ell}^t$ represents a case observed at $t$ but occurring at $t - d_{\ell}^t$. In what follows, we'll drop the superscript $t$ for simplicity. 

We assume that all delays follow a parametric distribution with cumulative distribution function (cdf):
\begin{equation}
d \sim G_D(\cdot|\theta_t).
\end{equation}
Furthermore, we assume the (unobserved) total number of cases at time $t$, $N_t$ (given by the sum over all  delays that lead to that time), has a parametric probability mass function (pmf): 
\begin{equation}
N_t \sim P_{N_t}(\cdot|\theta_t).
\end{equation}
The vector $\theta_t$ contains all of the model's parameters and its prior is $\pi(\theta_t)$. 

For the likelihood corresponding to a time $t$, we can average over all possible number of cases given that we have already observed $k$ delays:
\begin{equation}\label{likelihood_v1}
    \begin{aligned}
    \mathcal{L}_t(\theta_t  | d_1, \dots, d_ {k})  
    &\propto \pi(\theta_t)\cdot P(d_1, \dots, d_{k} | \theta_t)\\    
    & = \pi(\theta_t)\cdot \sum\limits_{n \geq k} P(d_1, \dots, d_{k} |\theta_t, N = n) \cdot P_N(n | \theta_t)   
    \end{aligned}
\end{equation}
When the number of cases, $n$, is fixed and using the fact that the first $k$ are less than $d^*$ while the last (unobserved) $n - k$ are larger, we obtain the following expression for the conditional probability:
\begin{equation}\label{gooddelay}
    \begin{aligned}
    P(d_1, \dots, d_{k} & |\theta_t, N_t = n)  
    \\ &= \left[ \prod\limits_{i=1}^k P\big(d_i - 1 < D_{(i)} \leq d_i|\theta_t\big)\right] \cdot \left[ \prod\limits_{i=k+1}^n P\big(D_{(i)} > d^{*}|\theta_t\big)\right] 
    \\ & = \binom{n}{k}\left[ \prod\limits_{i=1}^L \big[\Delta G_D(i|\theta_t)\big]^{m_i} \right] \cdot  \left[1 - G_D(d^*|\theta_t)\right]^{n-k},
    \end{aligned}
\end{equation}
where
\begin{equation}
    \Delta G_D(d_i|\theta_t) = \begin{cases}
        G_D(d_i|\theta_t) - G_D\big(d_i - 1|\theta_t) & \text{ if } G_{D} \text{ is continuous},\\\\
        g_D(d_i|\theta_t)  & \text{ if } G_{D} \text{ is discrete with mass } g_{D}.
    \end{cases}
\end{equation}
Integrating into \eqref{likelihood_v1} we obtain:
\begin{equation}\label{eq2}
    \begin{aligned}
    \mathcal{L}_t(\theta_t & | d_1, \dots, d_{k})  \propto   \pi(\theta_t) \cdot\left[ \prod\limits_{i=1}^L \big[\Delta G_D(i|\theta_t)\big]^{m_i} \right]\cdot  S_k(\theta_t)
    \end{aligned}
\end{equation}
with $m_i = \sum_{l=1}^k \mathbb{I}_{\{d_l = i\}}$ representing the total number of observed delays with value $i$, and the latent process, $S_k(\theta_t)$, is given by the following expression:
\begin{equation}\label{snk}
    S_k(\theta_t) = \sum\limits_{n \geq k} \binom{n}{k}   \cdot \left[1 - G_D(d^*|\theta_t)\right]^{n - k} \cdot P(N_t = n|\theta_t).
\end{equation}
where the total number of observed delays is $\sum_{i=1}^L m_i = k$. 

The log-likelihood for when $m_0$ cases have been observed with delay $0$, $m_1$ with delay $1$, and $m_L$ with delay $L$ is then:
\begin{equation}\label{loglik}
\begin{aligned}
    \ell_t(\theta_t & | d_1, \dots, d_{k}) = 
    \ln{\pi(\theta_t)} + \sum\limits_{l=0}^L m_l\cdot \ln{\left(\Delta G_D(d_l|\theta_t)\right)} + \ln{\left(S_k(\theta_t)\right)}.
\end{aligned}
\end{equation}
The expression for $S_k$ can be further simplified if $N_t$ has a $\text{Poisson}(\lambda_t)$ distribution:
\begin{equation}
    \ln S_k(\theta_t) = 
k \ln \lambda_t - \ln k! - G_D(d^* | \theta_t) \lambda_t
\end{equation}
or a $\text{NegativeBinomial}(r_t,p_t)$ distribution:
\begin{equation}
\begin{aligned}
\ln S_k(\theta_t) & = 
k \cdot \ln (1 - p_t) + r_t \cdot \ln p_t + \ln \binom{k + r_t - 1}{k} 
\\ & \quad - (k + r_t)\cdot \ln \big( p_t + G_D(d^*|\theta_t) \cdot (1 - p_t) \big).
\end{aligned}
\end{equation}

We remark that the likelihood in \eqref{eq2} can be interpreted as arising from a partially observed marked point process, where event times generate marks corresponding to reporting-delays. Under this view, inference is performed directly on the distribution of marks, rather than on the full event-report contingency table. This shift in perspective is central to our approach and underlies the extensions developed in later sections.

All derivations are shown in the appendix. 


\subsubsection{Missing report dates}

In general we consider that $L_1$ delay counts that have been accurately measured to the nearest integer: $m_1, m_2, \dots, m_{k_1}$ and $L_2$ delay counts that have been right-censored: $m_1^*, m_2^*, \dots, m_{k_2}^*$. Specifically, for each $m_j^*$, we consider that the only known information is that the delay is at most $j$ with $j \leq d^*$.   

In this scenario, we can rewrite \eqref{loglik} as:
\begin{equation}\label{logliksimplified}
    \begin{aligned}
        \ell(\theta_t & | m_0, \dots, m_{t}, m_0^{*}, \dots, m_{t}^{*}) 
        \\ & = \ln{\pi(\theta_t)} + \sum\limits_{i=1}^{L_1} m_i \cdot \ln{\left(\Delta G_D(i|\theta_t)\right)} +\sum\limits_{j=1}^{L_2} m_j^{*} \cdot \ln{\left(G_D(j|\theta_t)\right)} + \ln{\left(S_k(\theta_t)\right)}
    \end{aligned}
\end{equation}
which is equivalent to adding a censored term $P\big(D_{(i)} \leq d_i|\theta_t\big)$ to \eqref{gooddelay}. 

We remark that if the specific report date might be missing, the day that the nowcasting analysis is performed can be used as an upper bound to calculate the delay. 

\subsubsection{Parametric reporting-delay processes}

We model the reporting-delay distribution $G_D(\cdot \mid \theta_t)$ using a parametric family, which facilitates prior specification and the incorporation of structured temporal effects. Typical choices for $G_D$ include log-normal and generalized-gamma distributions, depending on the empirical characteristics of the delay distribution (e.g., skewness and tail behavior).

To capture temporal variation in reporting behavior, we parameterize the mean delay $\nu_t = \mathbb{E}[D \mid \theta_t]$ on the log scale:
\begin{equation}
\begin{aligned}
    \ln \nu_t &= v_0 + v^{\text{seasonal}}_{h(t)}.    
\end{aligned}
\end{equation}

Here, $v_0$ is a baseline intercept and $v_{h(t)}^{\text{seasonal}}$ represents seasonal effects. The function $h(t)$ indexes periodic effects such as day-of-the-week. For example, with weekly seasonality, $h(t) = \left[(t - 1) \bmod 7 \right] + 1$. Identifiability is ensured by imposing a sum-to-zero constraint on $\{v_{h(t)}^{\text{seasonal}}\}$. At the beginning of an epidemic, a constant $\nu_t$ can be used to stabilize inference.

We remark that while we parameterize the delay distribution through its mean, the full parameter vector $\theta_t$ may include additional shape or dispersion parameters.


\subsubsection{Nonparametric reporting-delay processes}

As an alternative to parametric specifications, we consider a flexible nonparametric model for the delay distribution based on a discrete probability mass function over a finite support. Let $L$ denote a maximum delay such that delays $d \in \{0,1,\dots,L\}$ are explicitly modeled. We place a Dirichlet prior on the delay probabilities:
\begin{equation}
    g_D = \big(g_D(0), \dots, g_D(L)\big) \sim \text{Dirichlet}(\alpha_0, \dots, \alpha_L).
\end{equation}

This formulation is widely used in nowcasting applications {\color{red} REF} due to its simplicity and ability to capture arbitrary delay shapes without strong parametric assumptions. However, it implicitly truncates the delay distribution at $L$, which complicates inference when late reports occur.

To address this limitation, we extend the support by introducing an additional category representing delays exceeding $L$. Specifically, we define an augmented probability vector:
\begin{equation}
    \tilde{g}_D = \big(g_D(0), \dots, g_D(L), g_D(L+1)\big) \sim \text{Dirichlet}(\alpha_0, \dots, \alpha_{L+1}),
\end{equation}
where $g_D(L+1)$ captures the total probability mass of delays greater than $L$. Conditional on belonging to this tail category, we model the excess delay using a continuous distribution. In this work, we adopt an exponential distribution with rate parameter 1, yielding a simple geometric-like tail behavior. The resulting cumulative distribution function is:
\begin{equation}
    \tilde{G}_D(x) =
    \begin{cases}
        \sum_{k=0}^{\lfloor x \rfloor} \tilde{g}_D(k), & \text{if } \lfloor x \rfloor \leq L, \\
        \sum_{k=0}^{L} \tilde{g}_D(k) + \tilde{g}_D(L+1)\left(1 - e^{-(x - (L+1))}\right), & \text{if } x > L.
    \end{cases}
\end{equation}

This construction preserves the flexibility of the discrete Dirichlet model on the observed range while handling right-censoring and late reports. It also enables the computation of probabilities for delays beyond the observed support, which is essential for likelihood-based inference in the presence of censoring.

\subsection{Epidemic process}

The latent epidemic process $\{N_t\}_t$ is modeled independently through a parametric count distribution $P(N_t \mid \theta_t)$, with mean $\mu_t = \mathbb{E}[N_t \mid \theta_t]$. Common choices include the Poisson and Negative Binomial distributions, the latter allowing for overdispersion.

We model the mean structure using a regression framework:
\begin{equation}
\mu_t = \gamma_0 + f(t) + \mathbf{X}_t^\top \boldsymbol{\gamma},
\end{equation}
where $\gamma_0$ is an intercept, $\mathbf{X}_t$ is a vector of observed covariates, and $\boldsymbol{\gamma}$ is a vector of regression coefficients with prior:
\begin{equation}
\boldsymbol{\gamma} \sim \mathcal{N}(\mathbf{0}, I_p).
\end{equation}

The function $f(t)$ captures temporal dependence in incidence and can be specified in several ways:

\begin{itemize}
    \item \textbf{Gaussian process (GP):}
    \begin{equation}
        f(t) \sim \mathcal{GP}(0, k(t,t')),
    \end{equation}
    allowing for flexible, nonparametric temporal structure.

    \item \textbf{Autoregressive trend:}
    \begin{equation}
    \begin{aligned}
        f(t) &= m_t, \\
        m_t &= \alpha\cdot  m_{t-1} + \epsilon_t, \quad \epsilon_t \sim \mathcal{N}(0, \sigma^2),
    \end{aligned}
    \end{equation}
    which captures smooth temporal evolution with limited complexity.

    \item \textbf{Mechanistic models (e.g., SIR-type):} In settings where epidemiological structure is important, $f(t)$ may be derived from compartmental models such as a Susceptible-Infected-Recovered framework. For example, incidence can be linked to transmission dynamics via:
    \begin{equation}
        f(t) \propto \beta S_t I_t,
    \end{equation}
    where $(S_t, I_t)$ evolve according to a discrete-time approximation of the underlying compartmental system.
\end{itemize}

This modular specification allows the epidemic process to range from purely statistical to mechanistic, depending on data availability and modeling goals.

\subsubsection{Handling strata}

In many applications, surveillance data are stratified by characteristics such as geography, age group, or reporting source. Let $s \in \{1, \dots, S\}$ index strata. We extend the model by defining stratum-specific incidence processes $N_t^{(s)}$ with means $\mu_t^{(s)}$:
\begin{equation}
    \mu_t^{(s)} = \gamma_0^{(s)} + f^{(s)}(t) + \mathbf{X}_t^{(s)\top} \boldsymbol{\gamma}^{(s)}.
\end{equation}

This formulation allows for heterogeneity across strata through separate regression coefficients and temporal components. Hierarchical priors may be introduced to share information across strata when appropriate.

For mechanistic models, stratification can be incorporated through coupled transmission dynamics. For example, in a multi-group SIR framework, the force of infection in stratum $s$ may depend on the total infectious population:
\begin{equation}
    \beta^{(s)} S_t^{(s)} \sum_{s'} I_t^{(s')},
\end{equation}
which captures cross-group interactions.

Overall, stratification is handled by extending both the mean structure and, when necessary, the dependence structure of the epidemic process, while retaining the same likelihood-based framework.

{\color{red} IGNORE THE PART IN RED
\subsubsection{Anomaly detection}

\paragraph{Scenario 1: new delay.}
Let $m_{\text{new}}$
denote the count of cases at event time $s \leq t$ that are newly reported with
this delay. By the Poisson thinning property, 
\begin{equation}\label{eq:s1_dist}
  M_{\mathrm{new}} \mid \lambda_s, \theta
  \;\sim\; \mathrm{Poisson}\!\bigl(\lambda_s \cdot \Delta G_D(d_{\mathrm{new}} \mid \theta)\bigr),
\end{equation}
%


\paragraph{Scenario 2: new time.}
At time $t+1$, we consider $m_{\mathrm{new}}$ and consider the posterior predictive distribution at $t+ 1$: 
%
\begin{equation}\label{eq:s2_dist}
  M_{\mathrm{new},s2} \mid \{\lambda_{t+1}^{(s)}\}, \theta
  \;\sim\; \mathrm{Poisson}\!\!\left(\Delta G_D(1 \mid \theta)
  \cdot \sum_{s=1}^{S} \lambda_{t+1}^{(s)}\right),?????
\end{equation}
%


\paragraph{Surprise scores.}

%
\begin{equation}
  \mathrm{LPD} = \log \mathbb{E}_\theta\bigl[p(m_{\mathrm{obs}} \mid \mathrm{rate}(\theta))\bigr]
  \approx \log\!\left(\frac{1}{S}\sum_{s=1}^{S} p\bigl(m_{\mathrm{obs}} \mid \mathrm{rate}^{(s)}\bigr)\right),
\end{equation}
%

\emph{posterior predictive $p$-values} 
%
\begin{equation}
  \mathrm{PPP}_{\mathrm{right}} = \mathbb{E}_\theta\bigl[P(M_{\mathrm{rep}} \geq m_{\mathrm{obs}} \mid \mathrm{rate}(\theta))\bigr],
  \qquad
  \mathrm{PPP}_{\mathrm{left}}  = \mathbb{E}_\theta\bigl[P(M_{\mathrm{rep}} \leq m_{\mathrm{obs}} \mid \mathrm{rate}(\theta))\bigr].
\end{equation}

The \emph{relative surprise} $\rho \in (0, 1]$:
%
\begin{equation}
  \rho = \exp\!\bigl(\mathrm{LPD}(m_{\mathrm{obs}}) - \mathrm{LPD}(m_{\mathrm{mode}})\bigr),
  \qquad
  m_{\mathrm{mode}} = \lfloor \mathbb{E}_\theta[\mathrm{rate}(\theta)] \rfloor.
\end{equation}
%

All code was developed in \texttt{R} {\color{red} (VERSION)} and is available in {\color{red} (REPO)}. 


}

\begin{comment} %version from march 19 2026
\subsubsection{Parametric reporting-delay processes}

For estimating the reporting-delay process $G_D$ we present several alternatives. In this section we focus on parametric processess which allow for inference when temporal effects such as day-of-the-week are included or when delays shift through time. Examples of $G_D$ include the lognormal, Weibull, and Gamma.  

Let $\nu_t$ denote the mean of the reporting-delay process. In general, we can write the mean of the delay as a process with a trend and temporal effects:
\begin{equation}\label{nudelay}
\begin{aligned}
    \ln \nu_t & = v_0 + v_{t}^{\text{trend}} + \beta^{\text{t-effect}}_{h(t)}  \\
    v_{t}^{\text{trend}} & = \phi \cdot v_{t-1}^{\text{trend}} + \epsilon_t.
\end{aligned}
\end{equation}  
where $\epsilon_t$ represents random noise, $v_{t}^{\text{trend}}$ a random walk that accomodates for time variations, and $\beta^{\text{t-effect}}_{h(t)}$ some coefficients for temporal effects. As an example, in the case of day-of-the-week, $h(t) = t \mod 7$ and the model has $\beta_1, \dots, \beta_7$ where $\beta_7 = -\sum_{i=1}^{6} \beta_i$ for identifiability. 

We note that the model described in \eqref{nudelay} can be constructed by parts (for example just including the intercept and the temporal effect). This is useful for contexts of low information where the user can intuitively parametrize the priors of a parametric model, particularly considering a constant delay. 

\subsubsection{Nonparametric reporting-delay processes}

Several nowcasting models such as {\color{red} REF SAM / NOBBS} consider a non-parametric, discrete, time-constant reporting-delay process. Specifically they consider that for a maximum delay of $L$, the delay pmf is given by:
\begin{equation}
g_D \sim \text{Dirichlet}(\alpha_L)
\end{equation}
where $g_D$ is a pmf over $\{0, \dots , L\}$. 

Here we extend that process to be able to compute censored probabilities (i.e. probabilities of delays larger than $L$) by considering an extended model of size $L + 1$:
\begin{equation}
\tilde{g}_D \sim \text{Dirichlet}(\alpha_{L+1})
\end{equation}
where the unobserved values are assigned an $\text{Exponential}(1)$ distribution. {\color{red} This is a geometric-rewrite!} This results in the following cdf:
\begin{equation}
    \tilde{G}_D( x ) = \begin{cases}
        \sum^{\lfloor x \rfloor}_{k=0} \tilde{g}_D(k) & \text{if } \lfloor x \rfloor \leq L \\
        \sum^{L}_{k=0} \tilde{g}_D(k) + 
        \big(1 - e^{-(x - (L + 1))}\big) \cdot \tilde{g}_D(L + 1) & \text{if } \lfloor x \rfloor > L, \\
    \end{cases}
\end{equation}
where the model corresponds to almost the same non-parametric prior for the delays less than $L$ and a parametric model that assigns an exponential probability once conditioned on arriving after the maximum delay. This allows to compute report dates greater than the maximum and to ask about the probability of a specific late arrival. 

\subsection{Epidemic process}

As previously discussed, the epidemic process can be modeled independently through $P(N_t|\theta_t)$. Specifically, we can use any function $f$ for describing the mean of the distribution ($\mu_t = \lambda_t$ in the Poisson case, $\mu_t = r_t (1 - p_t) / p_t$ for the Negative Binomial):
\begin{equation}
\begin{aligned}
\mu_t & = \gamma_0 + f(t) + \mathbf{X}_t^\top \boldsymbol{\gamma},
\end{aligned}
\end{equation}
where $\mathbf{X}_t$ is a vector of observed covariates at time $t$. 
The regression coefficients for the covariates can be modeled with standard normal priors:
\begin{equation}
\boldsymbol{\gamma} \sim \mathcal{N}(\mathbf{0}, I_p),
\end{equation}
where $p$ is the number of covariates and $I_p$ denotes the $p \times p$ identity matrix.
Some examples of functions $f(t)$ include:
\begin{itemize}
    \item A Gaussian process: 
    \begin{equation}
    f(t) \sim \mathcal{GP}\big(0, k(t,t')\big),
    \end{equation}
    with $k(t,t')$ its positive-definite covariance kernel.
    \item A model with an autoregresive trend:
    \begin{equation}
    \begin{aligned}
    f(t) & =  m_{t}\\
    m_{t} & = \alpha \cdot m_{t-1} + \epsilon_t
    \end{aligned}
    \end{equation}
    with $\epsilon_t$ a random error and $\alpha$ a parameter.
    \item Any SIR-like model:
    \begin{equation}
    f(t) = - \beta \cdot I_{t} \cdot S_{t}
    \end{equation}
    where the susceptibles and infected ($S_t, I_t$) are given by the system of equations:
    \begin{equation}
    \begin{aligned}
        S_t & = - \beta \cdot I_{t-1} \cdot S_{t-1}\\
        I_t & = \beta \cdot I_{t-1} \cdot S_{t-1} - \gamma \cdot I_{t-1}         
    \end{aligned}
    \end{equation}
\end{itemize}

\subsubsection{Handling strata}

Strata can be added for the epidemic process by computing stratified mean vectors $\mu^s_t$ that use as many parameters $\gamma_0, \gamma, \alpha. \epsilon$ as strata. A stratified Gaussian process can also be computed. As for the SIR model it can be changed by changing the infection rate $-\beta\cdot  I_{t-1} S_{t-1}$ to involve the sum of all infectious $-\beta^s \cdot S_{t-1}^s \cdot \sum_s I_{t-1}^s $ 
\end{comment}

\section{Simulation study}

\subsection{Model fit}
We performed a simulation study with a continuous time SIR model to represent a one-peak epidemic outbreak during a 125 day period. We spread across different reporting delays the number of cases with a Weibull$(2.70, 5.62)$ distribution (approximate mean of 5 days). To discretize into days, delays were rounded via the floor function. We compared three different approaches for the Poisson and Negative Binomial (NB) epidemic-processes: an autorregresive process of order 1, AR(1), a discrete Susceptible-Infected-Recovered (SIR) model, and a Hilbert Space Gaussian Process (HSGP). Furthermore we experimented with three different delay-process distributions: Lognormal, Generalized-Gamma, and Dirichlet (non-parametric). The model was fitted for each time $t = 6\dots, 125$ using the corresponding data observed up until that time. Fitting was done via a Laplace approximation {\color{red} REF}. We assessed the Weighted Interval Scoring (WIS) and Absolute Error (AE) of the fit by comparing the nowcasted estimates at each time against the simulated  "truth". 

\begin{figure}[!htb]
    \centering    \includegraphics[width=\linewidth]{images/plot_final_sims_complete_with_wis.pdf}
    \caption{Nowcasted values at notification time (colour) with 95\% credible interval versus final truth (points) and weighted interval score (WIS) for a simulated discrete SIR model with infection rate $\beta = 0.2$, recovery rate $\gamma = 0.1$, and overall population of $N = 1,000$ individuals. Delays were simulated using a Weibull$(2.70, 5.62)$ distribution with an approximate mean of $5$ days.}
    \label{fig:fitted_sim}
\end{figure}

Figure \ref{fig:fitted_sim} shows the fitted values and their 95\% credible intervals against the simulated "truth". It also features the WIS for each time-point. For all models, WIS was at its highest around the beginning of the epidemic peak. In general, HGSP outperformed the AR(1) and SIR models with an average $\text{WIS}_{\text{HGSP}}$ = 4.97 for the NB-Lognormal ($\text{WIS}_{\text{AR(1)}}$ = 7.46; $\text{WIS}_{\text{SIR}}$ = 9.22), 4.69 for the NB-Generalized-Gamma ($\text{WIS}_{\text{AR(1)}}$ = 9.68; $\text{WIS}_{\text{SIR}}$ = 7.51), and 4.68 for the NB-Dirichlet ($\text{WIS}_{\text{AR(1)}}$ = 8.77; $\text{WIS}_{\text{SIR}}$ = 8.07). In general, the NB likelihood yielded better scores than the Poisson model with 5.69 , 5.65, and 5.72 being the average $\text{WIS}_{\text{HGSP}}$ for the Lognormal, Generalized-Gamma, and Dirichlet under the Poisson model. The mean absolute error followed the same pattern. Table \ref{tab:wissim} describes these values as well as the decomposition into over and under dispersion, and the interval coverage. 
\begin{table}[!htb]
\centering
\caption{Nowcast comparison across different reporting-delay processes, epidemic-processes, and epidemic likelihood specifications for a simulated epidemic outbreak.}
\tiny
\setlength{\tabcolsep}{3pt}
\begin{tabular}{>{\raggedleft\arraybackslash}p{0.45in}|r | >{\raggedleft\arraybackslash}p{0.85in} | r r r r r r r r}
\toprule
\bf Epidemic process & \bf Likelihood & \bf Reporting-delay process  & \bf WIS & \bf Over & \bf Under & \bf Disp & \bf Bias & \bf IC 50\% & \bf IC 90\% & \bf AE \\
\midrule

\multirow{7}{*}{AR(1)}
& \multirow{3}{*}{NB}
& \cellcolor{delayC!25}Dirichlet & \cellcolor{delayC!25}9.57 & \cellcolor{delayC!25}0.43 & \cellcolor{delayC!25}7.54 & \cellcolor{delayC!25}1.60 & \cellcolor{delayC!25}-0.57 & \cellcolor{delayC!25}0.26 & \cellcolor{delayC!25}0.51 & \cellcolor{delayC!25}16.12 \\
& & \cellcolor{delayB!25}Gen-Gamma & \cellcolor{delayB!25}9.68 & \cellcolor{delayB!25}0.45 & \cellcolor{delayB!25}7.59 & \cellcolor{delayB!25}1.64 & \cellcolor{delayB!25}-0.58 & \cellcolor{delayB!25}0.23 & \cellcolor{delayB!25}0.50 & \cellcolor{delayB!25}16.30 \\
& & \cellcolor{delayA!25}Lognormal & \cellcolor{delayA!25}\bf  7.46 & \cellcolor{delayA!25}1.26 & \cellcolor{delayA!25}4.44 & \cellcolor{delayA!25}1.76 & \cellcolor{delayA!25}-0.36 & \cellcolor{delayA!25}0.35 & \cellcolor{delayA!25}0.61 & \cellcolor{delayA!25}13.01 \\

\\

& \multirow{3}{*}{Poisson}
& \cellcolor{delayC!25}Dirichlet & \cellcolor{delayC!25}\bf  8.77 & \cellcolor{delayC!25}0.45 & \cellcolor{delayC!25}2.35 & \cellcolor{delayC!25}5.97 & \cellcolor{delayC!25}-0.31 & \cellcolor{delayC!25}0.68 & \cellcolor{delayC!25}0.81 & \cellcolor{delayC!25}12.93 \\
& & \cellcolor{delayB!25}Gen-Gamma & \cellcolor{delayB!25}8.95 & \cellcolor{delayB!25}0.46 & \cellcolor{delayB!25}2.36 & \cellcolor{delayB!25}6.13 & \cellcolor{delayB!25}-0.30 & \cellcolor{delayB!25}0.66 & \cellcolor{delayB!25}0.81 & \cellcolor{delayB!25}13.08 \\
& & \cellcolor{delayA!25}Lognormal & \cellcolor{delayA!25}9.71 & \cellcolor{delayA!25}0.99 & \cellcolor{delayA!25}1.75 & \cellcolor{delayA!25}6.96 & \cellcolor{delayA!25}-0.13 & \cellcolor{delayA!25}0.68 & \cellcolor{delayA!25}0.82 & \cellcolor{delayA!25}13.03 \\

\midrule

\multirow{7}{*}{HSGP}
& \multirow{3}{*}{NB}
& \cellcolor{delayC!25}Dirichlet & \cellcolor{delayC!25}\bf  4.68 & \cellcolor{delayC!25}1.72 & \cellcolor{delayC!25}1.00 & \cellcolor{delayC!25}1.96 & \cellcolor{delayC!25}-0.02 & \cellcolor{delayC!25}0.40 & \cellcolor{delayC!25}0.82 & \cellcolor{delayC!25}10.10 \\
& & \cellcolor{delayB!25}Gen-Gamma & \cellcolor{delayB!25}4.69 & \cellcolor{delayB!25}1.67 & \cellcolor{delayB!25}1.08 & \cellcolor{delayB!25}1.94 & \cellcolor{delayB!25}-0.02 & \cellcolor{delayB!25}0.40 & \cellcolor{delayB!25}0.79 & \cellcolor{delayB!25}10.08 \\
& & \cellcolor{delayA!25}Lognormal & \cellcolor{delayA!25}4.97 & \cellcolor{delayA!25}2.12 & \cellcolor{delayA!25}0.86 & \cellcolor{delayA!25}1.99 & \cellcolor{delayA!25}0.05 & \cellcolor{delayA!25}0.41 & \cellcolor{delayA!25}0.81 & \cellcolor{delayA!25}10.56 \\

\\

& \multirow{3}{*}{Poisson}
& \cellcolor{delayC!25}Dirichlet & \cellcolor{delayC!25}5.72 & \cellcolor{delayC!25}0.49 & \cellcolor{delayC!25}1.83 & \cellcolor{delayC!25}3.40 & \cellcolor{delayC!25}-0.27 & \cellcolor{delayC!25}0.52 & \cellcolor{delayC!25}0.91 & \cellcolor{delayC!25}11.56 \\
& & \cellcolor{delayB!25}Gen-Gamma & \cellcolor{delayB!25}\bf 5.65 & \cellcolor{delayB!25}0.52 & \cellcolor{delayB!25}1.70 & \cellcolor{delayB!25}3.43 & \cellcolor{delayB!25}-0.25 & \cellcolor{delayB!25}0.52 & \cellcolor{delayB!25}0.92 & \cellcolor{delayB!25}11.31 \\
& & \cellcolor{delayA!25}Lognormal & \cellcolor{delayA!25}5.69 & \cellcolor{delayA!25}0.55 & \cellcolor{delayA!25}1.62 & \cellcolor{delayA!25}3.52 & \cellcolor{delayA!25}-0.22 & \cellcolor{delayA!25}0.55 & \cellcolor{delayA!25}0.93 & \cellcolor{delayA!25}11.42 \\

\midrule

\multirow{7}{*}{SIR}
& \multirow{3}{*}{NB}
& \cellcolor{delayC!25}Dirichlet & \cellcolor{delayC!25}7.82 & \cellcolor{delayC!25}0.28 & \cellcolor{delayC!25}5.87 & \cellcolor{delayC!25}1.68 & \cellcolor{delayC!25}-0.51 & \cellcolor{delayC!25}0.35 & \cellcolor{delayC!25}0.66 & \cellcolor{delayC!25}14.42 \\
& & \cellcolor{delayB!25}Gen-Gamma & \cellcolor{delayB!25}\bf 7.51 & \cellcolor{delayB!25}0.30 & \cellcolor{delayB!25}5.56 & \cellcolor{delayB!25}1.65 & \cellcolor{delayB!25}-0.49 & \cellcolor{delayB!25}0.36 & \cellcolor{delayB!25}0.65 & \cellcolor{delayB!25}13.98 \\
& & \cellcolor{delayA!25}Lognormal & \cellcolor{delayA!25}9.22 & \cellcolor{delayA!25}0.40 & \cellcolor{delayA!25}6.98 & \cellcolor{delayA!25}1.84 & \cellcolor{delayA!25}-0.39 & \cellcolor{delayA!25}0.33 & \cellcolor{delayA!25}0.64 & \cellcolor{delayA!25}15.43 \\

\\

& \multirow{3}{*}{Poisson}
& \cellcolor{delayC!25}Dirichlet & \cellcolor{delayC!25}\bf  8.07 & \cellcolor{delayC!25}0.19 & \cellcolor{delayC!25}3.01 & \cellcolor{delayC!25}4.87 & \cellcolor{delayC!25}-0.42 & \cellcolor{delayC!25}0.62 & \cellcolor{delayC!25}0.87 & \cellcolor{delayC!25}12.29 \\
& & \cellcolor{delayB!25}Gen-Gamma & \cellcolor{delayB!25}9.22 & \cellcolor{delayB!25}0.21 & \cellcolor{delayB!25}4.16 & \cellcolor{delayB!25}4.85 & \cellcolor{delayB!25}-0.39 & \cellcolor{delayB!25}0.61 & \cellcolor{delayB!25}0.86 & \cellcolor{delayB!25}14.00 \\
& & \cellcolor{delayA!25}Lognormal & \cellcolor{delayA!25}8.99 & \cellcolor{delayA!25}0.20 & \cellcolor{delayA!25}3.39 & \cellcolor{delayA!25}5.40 & \cellcolor{delayA!25}-0.33 & \cellcolor{delayA!25}0.69 & \cellcolor{delayA!25}0.86 & \cellcolor{delayA!25}14.48 \\

\bottomrule
\end{tabular}
\caption*{\tiny Reported values include Weighted Interval Scoring (WIS), Overdispersion (Over), Underdispersion (Under), Dispersion (Disp), Bias, Absolute Error (AE) and Interval Coverage (IC) for 50 and 90\% credible intervals. Models were fitted against a simulated discrete SIR model with infection rate $\beta = 0.2$, recovery rate $\gamma = 0.1$, and overall population of $N = 1,000$ individuals. Delays were simulated using a Weibull$(2.70, 5.62)$ distribution with an approximate mean of $5$ days.}
\label{tab:wissim}
\end{table}

{\color{red}
\subsection{Anomaly detection}
{\Large IGNORE THIS PART IN RED}
To assess the model's capabilities for anomaly detection we randomly selected $25$ times from the simulated SIR model and perturbed via two different options:
\begin{itemize}
    \item \textbf{Changed the average delay} by modeling it via a Weibull$(2.10, 2.26)$ with an approximate mean of 2.
    \item \textbf{Duplicating the number of cases} by substituting the case-counts at time $t + 1$ for delay $0$ from $m_0^{t+1}$ to $2m_0^{t+1}$.     
\end{itemize}
A model was fitted at time $t$ and then computed the relative surprise for the observations at time $t+1$. 
We classified an event as an anomaly if its relative surprise $\rho$ was less than $0.01$. 
}

\section{Real-Time Nowcasting Performance}

To assess our model we utilize the COVID-19 positive case database from Colombia's National Health Institute \cite{ins_colombia_covid_2026}. Data contains linelist information about each of the 6,390,971 COVID-19 cases diagnosed in the country from March 6th 2020 (first reported case) to January 16th 2024. It includes date of diagnosis (days), and date of first notification to the surveillance system (days). In general, suspected cases should be first notified to the surveillance system before the final diagnosis that results from testing the suspected cases. Hence the reporting delay is given by the difference between date of diagnosis and the notification date. 


\subsection{Model fit}
We fitted our models for each day with the information that would have been known by then (\textit{i.e} for each day $t$ we selected those with \texttt{diagnosis\_date $\leq t$ AND notification\_date $\leq t$}) and then compared against the "truth" defined as those with \texttt{diagnosis\_date $= t$} as notified by January 16th 2024 (the end of the data). For 242,080 cases where the diagnosis happened before the notification we assumed the diagnosis date was incorrect and set it as missing. Finally we evaluated the WIS of the predictions at each time $t$ for times $t$ (nowcast), $t-1,t-2, \dots, t-15$ (backcast). Figure \ref{fig:wis_colombia} describes the WIS for each of these times across different likelihood, reporting-delay and epidemic-process assumptions. As expected the WIS decreases as the lag increases (\textit{i.e.} with more information on the delay) except for the AR(1) models. The differences between the Poisson and the Negative Binomial tend to stabilize over time with the Poisson having consistently a larger WIS. 

\begin{figure}[!htb]
    \centering    \includegraphics[width=\linewidth]{images/wis_pt_COVID.pdf}
    \caption{Average weighted interval score for nowcasted values at notification time $t$ for time $t - \text{lag}$ for Colombia's COVID-19 positive case dataset with the reporting delay defined as the difference between \texttt{notification\_date} and \texttt{diagnosis\_date}.}
    \label{fig:wis_colombia}
\end{figure}

Table \ref{tab:wiscovid} describes the models' performance. Consistently with the simulation results, HGSP performed better in general. For example, under an NB likelihood and Generalized Gamma delay: ($\text{WIS}_{\text{HGSP}}$ = 1,670) than AR(1) and SIR epidemic processses ($\text{WIS}_{\text{AR(1)}}$ =  4,432; $\text{WIS}_{\text{SIR}}$ = 5,414). The only model competitive with the HSGP was the SIR model with a Dirichlet delay ($\text{WIS}_{\text{SIR}}$ = 1,745.23). The NB outperformed the Poisson likelihood in all cases. And in general the Dirichlet and Generalized Gamma delays fitted better than the Lognormal in the NB scenarios. The AR(1) had the most difficulty fitting the data and tended to generate  extreme credible intervals. Supplementary figure \ref{fig:fitted_sim_col} shows the nowcast's fit against the observed and simulated data. 

%NOTE-TO-SELF: RE-RUN all the models just to make sure these are the latest results
\begin{table}[!htb]
\centering
\caption{Nowcast comparison across different reporting-delay processes, epidemic-processes, and epidemic likelihood specifications for real-world epidemic outbreaks (Dengue, COVID-19, MPOX).}
\tiny
\setlength{\tabcolsep}{3pt}
\begin{tabular}{>{\raggedleft\arraybackslash}p{0.40in} | >{\raggedleft\arraybackslash}p{0.45in} | r | >{\raggedleft\arraybackslash}p{0.85in} | r r r r r r r r}
\toprule
\bf Disease & \bf Epidemic process & \bf Likelihood & \bf Reporting-delay process & \bf WIS & \bf Over & \bf Under & \bf Disp & \bf Bias & \bf IC 50\% & \bf IC 90\% & \bf AE \\
\midrule
 
multirow{23}{*}{Dengue}
 
& \multirow{7}{*}{AR(1)}
 
& \multirow{3}{*}{Poisson}
& \cellcolor{delayC!25}Dirichlet      & 74.40  & 23.89 & 23.32 & 27.20 & -0.76 & 0.19 & 0.32 & 105.73 \\
& & & \cellcolor{delayB!25}Gen-Gamma   & 38.88  & 0.02  & 22.34 & 16.52 & -0.86 & 0.13 & 0.46 & 34.89 \\
& & & \cellcolor{delayA!25}Lognormal   & 146.06 & 68.02 & 16.52 & 61.52 & -0.53 & 0.40 & \bf 0.63 & 214.30 \\
 
\\
 
& & \multirow{3}{*}{NB}
& \cellcolor{delayC!25}Dirichlet      & 34.09  & 0.00  & 33.87 & 0.22  & -1.00 & 0.00 & 0.00 & 35.78 \\
& & & \cellcolor{delayB!25}Gen-Gamma   & 35.89  & 0.00  & 33.23 & 2.66  & -0.97 & 0.03 & 0.04 & 35.30 \\
& & & \cellcolor{delayA!25}Lognormal   & 40.02  & 0.01  & 21.16 & 18.85 & -0.73 & 0.32 & 0.39 & \bf 30.50 \\
 
\midrule
 
& \multirow{7}{*}{HSGP}
 
& \multirow{3}{*}{Poisson}
& \cellcolor{delayC!25}Dirichlet      & 17.66  & 0.38  & 16.47 & 0.81  & -0.80 & 0.06 & 0.22 & 22.40 \\
& & & \cellcolor{delayB!25}Gen-Gamma   & 16.75  & 0.34  & 15.46 & 0.95  & -0.79 & 0.09 & 0.23 & 22.16 \\
& & & \cellcolor{delayA!25}Lognormal   & 17.11  & 0.40  & 15.86 & 0.84  & -0.80 & 0.06 & 0.23 & 21.93 \\
 
\\
 
& & \multirow{3}{*}{NB}
& \cellcolor{delayC!25}Dirichlet      & 14.28  & 0.11  & 12.17 & 2.00  & -0.64 & 0.26 & 0.46 & 22.01 \\
& & & \cellcolor{delayB!25}Gen-Gamma   & 13.87  & 0.13  & 11.69 & 2.05  & -0.63 & 0.26 & 0.44 & 21.56 \\
& & & \cellcolor{delayA!25}Lognormal   & \bf 13.34 & 0.16 & 11.05 & 2.13 & -0.62 & 0.27 & 0.49 & 21.50 \\
 
\midrule
 
& \multirow{7}{*}{SIR}
 
& \multirow{3}{*}{Poisson}
& \cellcolor{delayC!25}Dirichlet      & 27.78  & 0.08  & 26.93 & 0.77  & -0.88 & 0.02 & 0.23 & 34.56 \\
& & & \cellcolor{delayB!25}Gen-Gamma   & 23.15  & 0.08  & 21.56 & 1.51  & -0.81 & 0.11 & 0.36 & 33.10 \\
& & & \cellcolor{delayA!25}Lognormal   & 27.62  & 0.04  & 26.76 & 0.82  & -0.88 & 0.05 & 0.30 & 34.86 \\
 
\\
 
& & \multirow{3}{*}{NB}
& \cellcolor{delayC!25}Dirichlet      & 26.44  & 0.00  & 25.76 & 0.68  & -0.97 & 0.00 & 0.14 & 34.66 \\
& & & \cellcolor{delayB!25}Gen-Gamma   & 22.69  & 0.00  & 21.19 & 1.50  & -0.90 & 0.05 & 0.39 & 32.47 \\
& & & \cellcolor{delayA!25}Lognormal   & 24.63  & 0.00  & 23.79 & 0.84  & -0.95 & 0.01 & 0.23 & 33.92 \\
 
\midrule
 
& \multicolumn{2}{c|}{\cellcolor{nobbs!40}NobBS}
& \cellcolor{nobbs!40}Dirichlet       & \cellcolor{nobbs!40}33.87 & \cellcolor{nobbs!40}27.89 & \cellcolor{nobbs!40}0.00 & \cellcolor{nobbs!40}5.98 & \cellcolor{nobbs!40}0.89 & \cellcolor{nobbs!40}0.11 & \cellcolor{nobbs!40}0.11 & \cellcolor{nobbs!40}50.58 \\
 
& \multicolumn{2}{c|}{\cellcolor{epincast!40}Epinowcast (default)}
& \cellcolor{epincast!40}Lognormal     & \cellcolor{epincast!40}44.14 & \cellcolor{epincast!40}42.87 & \cellcolor{epincast!40}0.00 & \cellcolor{epincast!40}1.27 & \cellcolor{epincast!40}1.00 & \cellcolor{epincast!40}0.00 & \cellcolor{epincast!40}0.00 & \cellcolor{epincast!40}49.63 \\
 
& \multicolumn{2}{c|}{\cellcolor{epincast!40}Epinowcast (rw)}
& \cellcolor{epincast!40}Lognormal     & \cellcolor{epincast!40}44.62 & \cellcolor{epincast!40}43.24 & \cellcolor{epincast!40}0.00 & \cellcolor{epincast!40}1.37 & \cellcolor{epincast!40}1.00 & \cellcolor{epincast!40}0.00 & \cellcolor{epincast!40}0.00 & \cellcolor{epincast!40}51.84 \\
 
& \multicolumn{2}{c|}{\cellcolor{epincast!40}Epinowcast (point effect)}
& \cellcolor{epincast!40}Lognormal     & \cellcolor{epincast!40}45.26 & \cellcolor{epincast!40}44.19 & \cellcolor{epincast!40}0.00 & \cellcolor{epincast!40}1.07 & \cellcolor{epincast!40}1.00 & \cellcolor{epincast!40}0.00 & \cellcolor{epincast!40}0.00 & \cellcolor{epincast!40}51.26 \\
 
\midrule
\midrule
 
%% ============================================================
%% COVID
%% Best WIS:  NobBS / NB / Dirichlet = 1,043.92
%% Best AE:   AR1   / Poisson / Dirichlet = 2,094.02
%% Best IC90: HSGP  / NB / Dirichlet = 0.89
%% ============================================================
 
\multirow{23}{*}{COVID}
 
& \multirow{7}{*}{AR(1)}
 
& \multirow{3}{*}{Poisson}
& \cellcolor{delayC!25}Dirichlet      & 1{,}684.46 & 355.83 & 1{,}228.81 & 99.82  & -0.28 & 0.05 & 0.32 & \bf 2{,}094.02 \\
& & & \cellcolor{delayB!25}Gen-Gamma   & 1{,}798.39 & 395.17 & 1{,}305.52 & 97.69  & -0.35 & 0.18 & 0.31 & 2{,}175.54 \\
& & & \cellcolor{delayA!25}Lognormal   & 1{,}874.52 & 374.78 & 1{,}424.35 & 75.39  & -0.41 & 0.13 & 0.24 & 2{,}094.53 \\
 
\\
 
& & \multirow{3}{*}{NB}
& \cellcolor{delayC!25}Dirichlet      & 1{,}606.82 & 237.89 & 1{,}245.68 & 123.25 & -0.34 & 0.14 & 0.37 & 2{,}116.52 \\
& & & \cellcolor{delayB!25}Gen-Gamma   & 1{,}508.71 & 252.04 & 1{,}114.45 & 142.23 & -0.15 & 0.20 & 0.44 & 2{,}045.04 \\
& & & \cellcolor{delayA!25}Lognormal   & 1{,}620.47 & 269.39 & 1{,}210.44 & 140.64 & -0.32 & 0.20 & 0.33 & 2{,}170.88 \\
 
\midrule
 
& \multirow{7}{*}{HSGP}
 
& \multirow{3}{*}{Poisson}
& \cellcolor{delayC!25}Dirichlet      & 2{,}500.06 & 1{,}624.60 & 812.27  & 63.20  & 0.20  & 0.02 & 0.09 & 2{,}701.87 \\
& & & \cellcolor{delayB!25}Gen-Gamma   & 2{,}428.58 & 1{,}146.86 & 1{,}220.14 & 61.58 & -0.03 & 0.09 & 0.16 & 2{,}735.33 \\
& & & \cellcolor{delayA!25}Lognormal   & 2{,}319.69 & 1{,}092.66 & 1{,}205.85 & 21.18 & -0.19 & 0.02 & 0.07 & 2{,}425.04 \\
 
\\
 
& & \multirow{3}{*}{NB}
& \cellcolor{delayC!25}Dirichlet      & 1{,}121.29 & 371.53 & 144.35  & 605.41 & -0.03 & 0.55 & \bf 0.89 & 2{,}289.68 \\
& & & \cellcolor{delayB!25}Gen-Gamma   & 1{,}133.31 & 356.06 & 142.27  & 634.98 & 0.10  & 0.43 & 0.81 & 2{,}228.06 \\
& & & \cellcolor{delayA!25}Lognormal   & 1{,}376.34 & 523.70 & 259.92  & 592.72 & -0.06 & 0.36 & 0.84 & 3{,}044.54 \\
 
\midrule
 
& \multirow{7}{*}{SIR}
 
& \multirow{3}{*}{Poisson}
& \cellcolor{delayC!25}Dirichlet      & 2{,}447.53 & 92.41  & 2{,}040.62 & 314.49 & -0.68 & 0.18 & 0.42 & 3{,}530.66 \\
& & & \cellcolor{delayB!25}Gen-Gamma   & 2{,}584.43 & 495.93 & 2{,}035.67 & 52.82  & -0.42 & 0.06 & 0.17 & 2{,}833.04 \\
& & & \cellcolor{delayA!25}Lognormal   & 2{,}232.16 & 348.01 & 1{,}624.41 & 259.74 & -0.63 & 0.16 & 0.33 & 2{,}907.71 \\
 
\\
 
& & \multirow{3}{*}{NB}
& \cellcolor{delayC!25}Dirichlet      & 1{,}957.00 & 32.65  & 1{,}434.44 & 489.91 & -0.68 & 0.29 & 0.73 & 3{,}451.94 \\
& & & \cellcolor{delayB!25}Gen-Gamma   & 1{,}916.14 & 108.77 & 1{,}484.84 & 322.53 & -0.54 & 0.20 & 0.60 & 3{,}268.12 \\
& & & \cellcolor{delayA!25}Lognormal   & 2{,}188.50 & 177.30 & 1{,}696.41 & 314.80 & -0.58 & 0.21 & 0.52 & 3{,}514.23 \\
 
\midrule
 
& \multicolumn{2}{c|}{\cellcolor{nobbs!40}NobBS}
& \cellcolor{nobbs!40}Dirichlet       & \cellcolor{nobbs!40}\bf 1{,}043.92 & \cellcolor{nobbs!40}390.92 & \cellcolor{nobbs!40}69.27 & \cellcolor{nobbs!40}583.73 & \cellcolor{nobbs!40}0.29 & \cellcolor{nobbs!40}0.49 & \cellcolor{nobbs!40}0.87 & \cellcolor{nobbs!40}\bf 2{,}200.80 \\
 
& \multicolumn{2}{c|}{\cellcolor{epincast!40}Epinowcast (default)}
& \cellcolor{epincast!40}Lognormal     & \cellcolor{epincast!40}35{,}542.19 & \cellcolor{epincast!40}33{,}786.03 & \cellcolor{epincast!40}263.59 & \cellcolor{epincast!40}1{,}492.57 & \cellcolor{epincast!40}0.58 & \cellcolor{epincast!40}0.07 & \cellcolor{epincast!40}0.13 & \cellcolor{epincast!40}49{,}930.98 \\
 
& \multicolumn{2}{c|}{\cellcolor{epincast!40}Epinowcast (rw)}
& \cellcolor{epincast!40}Lognormal     & \cellcolor{epincast!40}38{,}194.70 & \cellcolor{epincast!40}31{,}859.06 & \cellcolor{epincast!40}40.53 & \cellcolor{epincast!40}6{,}295.11 & \cellcolor{epincast!40}0.80 & \cellcolor{epincast!40}0.09 & \cellcolor{epincast!40}0.12 & \cellcolor{epincast!40}57{,}593.93 \\
 
& \multicolumn{2}{c|}{\cellcolor{epincast!40}Epinowcast (point effect)}
& \cellcolor{epincast!40}Lognormal     & \cellcolor{epincast!40}26{,}745.88 & \cellcolor{epincast!40}14{,}175.74 & \cellcolor{epincast!40}123.90 & \cellcolor{epincast!40}12{,}446.24 & \cellcolor{epincast!40}0.83 & \cellcolor{epincast!40}0.05 & \cellcolor{epincast!40}0.22 & \cellcolor{epincast!40}32{,}235.32 \\
 
\midrule
\midrule
 
%% ============================================================
%% MPOX
%% Best WIS:  SIR / Poisson / Dirichlet = 8.43
%% Best AE:   SIR / Poisson / Dirichlet = 13.55
%% Best IC90: AR1 / Poisson / Gen-Gamma = 1.00 (tie with AR1/Poisson/Dirichlet)
%% ============================================================
 
\multirow{23}{*}{MPOX}
 
& \multirow{7}{*}{AR(1)}
 
& \multirow{3}{*}{Poisson}
& \cellcolor{delayC!25}Dirichlet      & 14.64 & 0.47 & 5.53  & 8.64  & -0.32 & 0.51 & \bf 1.00 & 24.13 \\
& & & \cellcolor{delayB!25}Gen-Gamma   & 19.24 & 0.42 & 7.30  & 11.53 & -0.24 & 0.47 & \bf 1.00 & 28.60 \\
& & & \cellcolor{delayA!25}Lognormal   & 15.14 & 0.46 & 6.38  & 8.30  & -0.35 & 0.43 & 0.97 & 25.39 \\
 
\\
 
& & \multirow{3}{*}{NB}
& \cellcolor{delayC!25}Dirichlet      & 19.22 & 0.55 & 17.42 & 1.25  & -0.50 & 0.27 & 0.41 & 24.39 \\
& & & \cellcolor{delayB!25}Gen-Gamma   & 24.44 & 0.56 & 22.58 & 1.30  & -0.46 & 0.27 & 0.44 & 29.76 \\
& & & \cellcolor{delayA!25}Lognormal   & 21.15 & 1.65 & 17.98 & 1.51  & -0.35 & 0.23 & 0.42 & 27.57 \\
 
\midrule
 
& \multirow{7}{*}{HSGP}
 
& \multirow{3}{*}{Poisson}
& \cellcolor{delayC!25}Dirichlet      & 14.79 & 0.35 & 12.62 & 1.82  & -0.30 & 0.35 & 0.69 & 21.61 \\
& & & \cellcolor{delayB!25}Gen-Gamma   & 15.97 & 0.25 & 14.08 & 1.64  & -0.39 & 0.30 & 0.67 & 22.56 \\
& & & \cellcolor{delayA!25}Lognormal   & 15.15 & 0.25 & 13.22 & 1.68  & -0.38 & 0.30 & 0.67 & 21.98 \\
 
\\
 
& & \multirow{3}{*}{NB}
& \cellcolor{delayC!25}Dirichlet      & 11.84 & 0.25 & 9.21  & 2.38  & -0.30 & 0.40 & 0.75 & 20.70 \\
& & & \cellcolor{delayB!25}Gen-Gamma   & 12.16 & 0.21 & 9.60  & 2.34  & -0.32 & 0.39 & 0.77 & 20.89 \\
& & & \cellcolor{delayA!25}Lognormal   & 11.89 & 0.20 & 9.32  & 2.37  & -0.34 & 0.41 & 0.77 & 20.76 \\
 
\midrule
 
& \multirow{7}{*}{SIR}
 
& \multirow{3}{*}{Poisson}
& \cellcolor{delayC!25}Dirichlet      & \bf 8.43 & 0.94 & 6.24  & 1.26  & -0.33 & 0.31 & 0.62 & \bf 13.55 \\
& & & \cellcolor{delayB!25}Gen-Gamma   & 9.12  & 0.76 & 7.09  & 1.27  & -0.36 & 0.27 & 0.61 & 14.58 \\
& & & \cellcolor{delayA!25}Lognormal   & 8.42  & 0.69 & 6.37  & 1.36  & -0.36 & 0.27 & 0.60 & 13.80 \\
 
\\
 
& & \multirow{3}{*}{NB}
& \cellcolor{delayC!25}Dirichlet      & 12.86 & 0.25 & 7.54  & 5.06  & -0.46 & 0.40 & 0.93 & 24.55 \\
& & & \cellcolor{delayB!25}Gen-Gamma   & 13.58 & 0.41 & 8.05  & 5.13  & -0.44 & 0.38 & 0.94 & 25.34 \\
& & & \cellcolor{delayA!25}Lognormal   & 13.30 & 0.40 & 7.48  & 5.42  & -0.42 & 0.35 & 0.94 & 25.03 \\
 
\midrule
 
& \multicolumn{2}{c|}{\cellcolor{nobbs!40}NobBS}
& \cellcolor{nobbs!40}Dirichlet       & \cellcolor{nobbs!40}29.98 & \cellcolor{nobbs!40}1.05 & \cellcolor{nobbs!40}27.55 & \cellcolor{nobbs!40}1.38 & \cellcolor{nobbs!40}-0.17 & \cellcolor{nobbs!40}0.20 & \cellcolor{nobbs!40}0.45 & \cellcolor{nobbs!40}33.15 \\
 
& \multicolumn{2}{c|}{\cellcolor{epincast!40}Epinowcast (default)}
& \cellcolor{epincast!40}Lognormal     & \cellcolor{epincast!40}25.13 & \cellcolor{epincast!40}15.58 & \cellcolor{epincast!40}3.10 & \cellcolor{epincast!40}6.46 & \cellcolor{epincast!40}0.39 & \cellcolor{epincast!40}0.18 & \cellcolor{epincast!40}0.45 & \cellcolor{epincast!40}40.62 \\
 
& \multicolumn{2}{c|}{\cellcolor{epincast!40}Epinowcast (rw)}
& \cellcolor{epincast!40}Lognormal     & \cellcolor{epincast!40}27.96 & \cellcolor{epincast!40}13.55 & \cellcolor{epincast!40}4.37 & \cellcolor{epincast!40}10.04 & \cellcolor{epincast!40}0.30 & \cellcolor{epincast!40}0.24 & \cellcolor{epincast!40}0.44 & \cellcolor{epincast!40}36.55 \\
 
& \multicolumn{2}{c|}{\cellcolor{epincast!40}Epinowcast (point effect)}
& \cellcolor{epincast!40}Lognormal     & \cellcolor{epincast!40}17.72 & \cellcolor{epincast!40}9.50 & \cellcolor{epincast!40}0.67 & \cellcolor{epincast!40}7.55 & \cellcolor{epincast!40}0.61 & \cellcolor{epincast!40}0.27 & \cellcolor{epincast!40}0.58 & \cellcolor{epincast!40}30.32 \\
 
\bottomrule
\end{tabular}
\caption*{\tiny Reported values include Weighted Interval Scoring (WIS), Overdispersion (Over), Underdispersion (Under), Dispersion (Disp), Bias, Absolute Error (AE) and Interval Coverage (IC) for 50 and 90\% credible intervals. Bold values indicate the best (lowest) WIS, best (lowest) AE, and best (highest) IC 90\% per disease. Epinowcast variants and NobBS are shown with distinct background shading.}
\label{tab:wisreal}
\end{table}


{\color{red}
\subsection{Comparison to other methods}

In addition to the predictions for our model we also utilized the \texttt{Nobbs} and \texttt{epinowcast} packages to generate predictions for the cases. Comparison was made against our XXX and YYY models. The full specification for the \texttt{Nobbs} and \texttt{epinowcast} calls can be found in the appendix. 

}
\section{Discussion}

{\color{red} Compare number of operations with nobbs

Produces smoothed (filtered) versions

Identifiability and priors

}




\begin{comment}
The epidemic process can also be any form of a dynamic SIR model:
\begin{equation}
\begin{aligned}
    BLAH
\end{aligned}    
\end{equation}
\end{comment}

\bibliographystyle{plain}
\bibliography{sample}

\clearpage % Ensures the appendix starts on a new page
\appendix
\section{Additional figures}

\begin{figure}[!htb]
    \centering    \includegraphics[width=\linewidth]{images/plot_final_sims_complete_with_wis_covid.pdf}
    \caption{Nowcasted values at notification time (colour) with 95\% credible interval versus final truth (points) and weighted interval score (WIS) for a simulated discrete SIR model with infection rate $\beta = 0.2$, recovery rate $\gamma = 0.1$, and overall population of $N = 1,000$ individuals. Delays were simulated using a Weibull$(2.70, 5.62)$ distribution with an approximate mean of $5$ days.}
    \label{fig:fitted_sim_col}
\end{figure}

\end{document}